\section{An actual-source reciprocal bound and a conditional interior-cluster obstruction}
\label{sec:full-source-cluster-obstruction}

\paragraph{Status.}
The complete author proof reading and two independent whole-proof
readings have passed.  Their exact source and final-scope check
records are kept separately.

This result concerns the original full nontrivial common-sector
map, not the completed ideal return or the whole-quotient return.
It extracts a quantitative necessary condition from a family of
genuine compact inputs.  The cluster conclusion below is conditional
on the stated zero configuration; no occurrence of that configuration
for a Hecke function is asserted.

Keep the number field, nontrivial primitive finite-order character
$\eta$, dual character, structural $S$, and actual source maps of
(MA-MO.550)--(MA-MO.552).  In particular $\phi$ is a fixed real
smooth step, equal to one sufficiently far left and zero sufficiently
far right, and $q=1-\phi$.  Write
\[
\begin{gathered}
 A_\omega(z)=E_{\omega,S}(1/2+z)L_K(1/2+z,\omega),
       \qquad \omega\in\{\eta,\bar\eta\},\\
 H=H_{\Sigma,a,s},\qquad
 w_{a,s}(x)=(e^{2ax}+e^{-2ax})(1+x^2)^{s/2},\\
 0<a<1/2,\qquad s\ge0,\qquad
 p=\frac{s+1}{2},\qquad
 0<\delta_0<\min\{1,(1/2-a)/2\},\\
 D_g^0=C_g(P_+-P_-):\mathcal E\longrightarrow\mathcal E,
       \qquad \mathcal D=C_c^\infty([1,2]).
\end{gathered}
\tag{MA-MO.1200}
\]
The meaning of $\mathcal D$ is the closed Fr\'echet space of
smooth functions on $\mathbb R$ supported in $[1,2]$.
Its real subspace will also be used.  Every test in this space
has $g(0)=0$.

Assume, for the comparison below, that every completed nontrivial
zero of both sectors lies in $|\Re z|\le a$ in centered coordinates.
The additional full-retention condition means that every zero on
an edge has multiplicity $m$ satisfying $m<p$.  The latter condition
is necessary by MA-MO-025, but is not used to assume any inverse
estimate near an edge.  Put $\nu(T)=2+|T|$ and define the condition
\[
 \begin{gathered}
 (\mathrm{BR}_{a,s,\eta}):\quad
 \exists C,N<\infty\quad
 \delta^p|A_\eta(a+\delta+iT)|^{-1}
       \le C\nu(T)^N\\
 \hfill(T\in\mathbb R,\quad0<\delta\le\delta_0).
 \end{gathered}
 \tag{MA-MO.1201}
\]
This is a two-variable estimate from the open right half-plane.
It does not assign a boundary reciprocal multiplier or a boundary
distribution to arbitrary Hilbert data.

\paragraph{MA-MO-033 (a full-source family requires boundary-uniform reciprocal control).}
Under the closed-strip premise above, if $D_g^0$ has a bounded
extension on $H$ for every complex $g\in\mathcal D$, then
$(\mathrm{BR}_{a,s,\eta})$ holds.  The same conclusion follows
if boundedness is assumed only for every real $g\in\mathcal D$.
The positive-support family controls the right edge of $\eta$.
By reflection, boundedness for the negative-support $\eta$ family
controls the right edge of $\bar\eta$, equivalently the left
edge of $\eta$.  No dual bound is inferred from the positive
family alone.

If (MA-MO.1201) fails, there is a fixed sequence
\[
 e_n\in C_c^\infty(\mathbb R)\subset\mathcal E,
           \qquad\|e_n\|_H=1,
 \tag{MA-MO.1202}
\]
such that the set of $g\in\mathcal D$ for which
$\sup_n\|D_g^0e_n\|_H=\infty$ is a dense $G_\delta$.
This holds separately in the real and complex test spaces.
Thus these fixed tests admit no bounded extension of their actual
source action.  The inputs in (MA-MO.1202) are normalized finite
translation-orbit integrals, specified in the proof.

A sufficient, explicitly hypothetical configuration for this failure
is a sequence $|T_n|\to\infty$, $d_n\downarrow0$, with at least
$k_n$ centered zeros, counted with multiplicity, in
\[
 \left\{z:\left|z-(a-d_n+iT_n)\right|\le d_n/2\right\},
 \qquad
 \frac{(k_n-p)\log(1/d_n)}{\log\nu(T_n)}\longrightarrow\infty.
 \tag{MA-MO.1203}
\]
All zeros in each displayed cluster are strictly inside the strip.
Here and in the proof discard a finite prefix so that
$d_n<\min\{\delta_0,4a/3\}$; this makes both edge inequalities
strict and permits the choice $\delta=d_n$.
In particular, a fixed integer $k>p$ and
$\log(1/d_n)/\log\nu(T_n)\to\infty$ suffice.
Unbounded $k_n$ with suitably small polynomial $d_n$ also suffice.
The configuration may consist of distinct simple zeros.  Neither
these possibilities nor the full-retention premise are asserted
to occur arithmetically.

\begin{proof}
\emph{1. Fixed compact supports and uniform noncancellation.}
Let $r=\sup\{x:\phi(x)\ne0\}$.  Choose a compact interval
$J$ in $(r-\log2,r)$ on which $\phi$ has a constant nonzero
sign and is bounded away from zero.  Such an interval exists
by the definition of $r$ and continuity.  Choose nonnegative,
nonzero $g_0\in C_c^\infty((1,2))$ and
$v_0\in C_c^\infty(\mathbb R)$, with sufficiently small
supports, so that
\[
 k_0=C_{g_0^\vee}v_0\ne0,\qquad
             k_0\ge0,\qquad \operatorname{supp}k_0\subset J.
 \tag{MA-MO.1204}
\]
For example, place the support of $g_0$ near $t_*$ in $(1,2)$
and that of $v_0$ near $x_*+t_*$, with $x_*$ in the interior
of $J$.  Their support difference then lies in $J$.
Fix also $0\ne h_0\ge0$ compact and smooth, with
$\operatorname{supp}h_0\subset[\alpha,\beta]$.
For real $T$ set
\[
 g_T(t)=e^{-iTt}g_0(t),\qquad
 v_T(x)=e^{iTx}v_0(x),\qquad
 h_T(x)=e^{-iTx}h_0(x).
 \tag{MA-MO.1205}
\]
The bilinear transpose in (MA-MO.557) gives exactly
$C_{g_T^\vee}v_T=e^{iTx}k_0(x)$.  The norm of the output
functional $\ell_{v_T}(y)=\int yv_T$ on $H$ is independent
of $T$.  Also
\[
 \|g_T\|_{C^j}\le C_j\nu(T)^j
                   \qquad(j\ge0).
 \tag{MA-MO.1206}
\]

Use the literal $F,W,H_h,V_g,B_v,\mathcal R_{h,v}$ of
(MA-MO.558)--(MA-MO.567), with these modulated tests.
Every nonunit ideal has norm at least two.  Since
$\operatorname{supp}k_0\subset(r-\log2,r)$, the only term
of $W$ meeting $\operatorname{supp}\phi$ is the unit ideal.
At $w=a+\delta+iT$ this proves
\[
 \begin{aligned}
 H_{h_T}(w)&=\int h_0(x)e^{(a+\delta)x}\,dx,\\
 B_{v_T}(w)&=\int_J\phi(x)k_0(x)e^{-(a+\delta)x}\,dx,\\
 V_{g_T}(-w)&=\int_Jk_0(x)e^{-(a+\delta)x}\,dx.
 \end{aligned}
 \tag{MA-MO.1207}
\]
Thus, for $0<\delta\le\delta_0$, both of the first two
absolute values are bounded below by fixed positive constants,
and all three are bounded above.  This is exact unit-ideal
noncancellation after modulation, including for complex $g_T$;
no estimate of a global inverse has been used.

Choose $t_0$ as in (MA-MO.561), using these fixed supports,
and enlarge it so that $\beta+t_0\ge1$.
The dual source term on $U_th_T$ is identically zero for
$t\ge t_0$, before applying the actual source Fourier map.
The direct convolution pairing also vanishes there.
The compact region (MA-MO.565) is independent of $T$.
On that region both ideal sums are finite, with uniformly
bounded absolute values: modulation only changes their terms
by phases.  The exponential factor in the remainder has
bounded modulus for $a\le\Re w\le a+\delta_0$.
Consequently
\[
       |\mathcal R_{h_T,v_T}(a+\delta+iT)|\le C
                    \quad(T\in\mathbb R,\ 0<\delta\le\delta_0).
 \tag{MA-MO.1208}
\]
This bound uses the compact double integral, not a bound on
individual partial-fraction coefficients.

\emph{2. The orbit identity on its actual source domain.}
Put $w=a+\delta+iT$ and, for $L>t_0$, define
\[
 \begin{aligned}
 u_{\delta,T,L}
      &=\int_{t_0}^{L}e^{-wt}U_th_T\,dt,\\
 u_{\delta,T}(x)
      &=e^{-wx}\int_{-\infty}^{x-t_0}
                         e^{(a+\delta)y}h_0(y)\,dy .
 \end{aligned}
 \tag{MA-MO.1209}
\]
The first vector is compact and smooth, supported in
$[\alpha+t_0,\beta+L]$.  It belongs to the original smooth
common sector, and (MA-MO.552) acts on it using the literal
wide inverse and actual local Fourier pair.
Changing variables gives its more explicit expression
\[
 u_{\delta,T,L}(x)=e^{-wx}
       \int_{x-L}^{x-t_0}e^{(a+\delta)y}h_0(y)\,dy.
 \tag{MA-MO.1210}
\]
Positivity of $h_0$ shows that its absolute value is bounded
by that of $u_{\delta,T}$.  The latter vanishes left of
$\alpha+t_0$ and has the exact tail
$H_{h_T}(w)e^{-wx}$ right of $\beta+t_0$.
The compact middle interval is uniformly bounded.  Retaining
both exponentials in the sum norm, exactly as in (MA-MO.614),
gives
\[
 \|u_{\delta,T,L}\|_H
       \le\|u_{\delta,T}\|_H\le C\delta^{-p}.
 \tag{MA-MO.1211}
\]
All constants here are independent of $T,\delta,L$.
The finite vectors converge to $u_{\delta,T}$ in $H$ by
dominated convergence.

There is also an independent justification for the scalar
output limit which does not assume the desired extension on
$H$.  For fixed $\delta>0$, take $b=a+\delta/2$.
The closed-strip hypothesis gives the strict right-half-plane
margin required by MA-MO-022 at width $b$, for both sectors.
That theorem supplies the actual source extension on
$H_{\Sigma,b,s}$, with constants allowed to depend on $b$.
The vectors in (MA-MO.1209) converge in this wider norm too,
since $\Re w>b$.  Equivalently, the actual paired function
$f(t)=\ell_{v_T}(D_{g_T}^0U_th_T)$ has an exponential bound
with every fixed exponent strictly above $a$.  Hence its
Laplace integral converges whenever $\Re w>a$.

The wide identity (MA-MO.567), initially valid for
$\Re w>1/2$, therefore continues holomorphically to
$\Re w>a$; $A_\eta$ has no zero there.  At the specified
$w$ it yields
\[
 \begin{aligned}
 \lim_{L\to\infty}\ell_{v_T}(D_{g_T}^0u_{\delta,T,L})
   &=\mathfrak F_T(w),\\
 \mathfrak F_T(w)
   &=H_{h_T}(w)V_{g_T}(-w)
       -\frac{H_{h_T}(w)B_{v_T}(w)}{A_\eta(w)}
       -\mathcal R_{h_T,v_T}(w),\\
 |\mathfrak F_T(w)|&\ge c|A_\eta(w)|^{-1}-C.
 \end{aligned}
 \tag{MA-MO.1212}
\]
Finite-interval integration is on the actual compact source.
The wider-space result is used only to justify this scalar
limit for each fixed $\delta$; no uniform $b\downarrow a$
operator estimate is imported.

\emph{3. Baire gives the required uniform test estimate.}
Suppose first that every complex test in $\mathcal D$ has
a bounded extension on $H$.  Uniqueness follows from density
of compact smooth inputs.  For each fixed compact input $u$,
the map $g\mapsto D_g^0u$ is continuous into $H$:
$D_g^0u=C_gD^0u$, with $D^0u\in\mathcal E$, and the
weighted translation bound controls convolution on a fixed
test-support interval.  Therefore the sets
\[
 \mathcal F_m=\{g\in\mathcal D:
      \|D_g^0u\|_H\le m\|u\|_H
                 \text{ for every }u\in C_c^\infty\}
 \tag{MA-MO.1213}
\]
are closed.  They cover $\mathcal D$ under the hypothesis.
Baire's theorem, subtraction of a point in an interior ball,
linearity, and scaling give some finite $N,C$ such that
\[
                \|D_g\|_{H\to H}\le C\|g\|_{C^N}
                              \qquad(g\in\mathcal D).
 \tag{MA-MO.1214}
\]
The identical argument in the real Fr\'echet space gives this
bound first for real tests; writing a complex test as its real
and imaginary parts extends it with at most a fixed factor.
Thus the real-only hypothesis suffices as well.

Apply the bound to (MA-MO.1212), then use (MA-MO.1211), the
constant norm of $\ell_{v_T}$, and (MA-MO.1206).  It follows that
\[
 c|A_\eta(a+\delta+iT)|^{-1}-C
      \le C'\nu(T)^N\delta^{-p}.
 \tag{MA-MO.1215}
\]
Multiplication by $\delta^p$, with $\delta\le\delta_0$,
proves (MA-MO.1201).

\emph{4. A fixed genuine-core sequence and residual bad tests.}
Assume that (MA-MO.1201) fails.  Choose pairs $(T_n,\delta_n)$
so that $\delta_n^p|A_\eta(a+\delta_n+iT_n)|^{-1}$ exceeds
any prescribed constant times $n\nu(T_n)^n$.
By (MA-MO.1212), choose a finite $L_n>t_0$ such that its
actual scalar output differs from $\mathfrak F_{T_n}$ by
at most one, and put
\[
 e_n=\frac{u_{\delta_n,T_n,L_n}}
              {\|u_{\delta_n,T_n,L_n}\|_H}.
 \tag{MA-MO.1216}
\]
The denominator is nonzero by (MA-MO.1210).  Increasing the
chosen reciprocal lower bounds if necessary, equations
(MA-MO.1211)--(MA-MO.1212) and the uniform functional norm give
\[
                  \|D_{g_{T_n}}^0e_n\|_H
                                    \ge n\nu(T_n)^n.
 \tag{MA-MO.1217}
\]
This construction uses only finite original-source inputs.
It may choose $L_n$ adaptively; no uniform truncation rate is claimed.

For these now fixed inputs define the closed sets
\[
 \mathcal B_m=\{g\in\mathcal D:
                    \sup_n\|D_g^0e_n\|_H\le m\}.
 \tag{MA-MO.1218}
\]
If one had interior, the same subtraction and scaling argument
would give $\sup_n\|D_g^0e_n\|_H\le C\|g\|_{C^N}$ for all
tests.  Applied to $g_{T_n}$, this contradicts
(MA-MO.1206) and (MA-MO.1217) for $n>N$.
Every $\mathcal B_m$ is therefore nowhere dense.
Their union is exactly the set of tests with bounded outputs
on the fixed sequence.  The complement is a dense $G_\delta$.
For real tests the same interior argument complexifies before
using $g_{T_n}$, so the real statement follows too.

If full individual retention is additionally assumed, reciprocal
growth failure cannot be confined to bounded heights: near a
boundary zero of order $m<p$, $\delta^p|A(a+\delta+iT)|^{-1}$
is bounded locally, and away from finitely many such zeros
compactness applies.  Thus in that setting the pairs may be
chosen with $|T_n|\to\infty$.  This observation is not needed
for the genuine-core construction itself.

\emph{5. A conditional cluster forces failure.}
We verify the finite-disk estimate needed here directly.
At large $|T|$, center the uncentered finite-$L$ variable at
$2+iT$ and choose a zero-free boundary radius between $12$
and $13$.  Divide $L_K$ by the finite disk Blaschke product
of its zeros, with multiplicity.  The resulting holomorphic
factor $H_T$ has polynomial upper bound $C\nu(T)^D$ on
the disk, by the finite-$L$ polynomial strip bound and the
maximum principle; each Blaschke factor has modulus one
on its boundary.  On the inner half disk each normalized
factor belonging to a zero $\rho$ satisfies
\[
      |B_{\rho,T}(w)|\le C|w-\rho|,
                       \qquad |B_{\rho,T}(w)|\le1.
 \tag{MA-MO.1219}
\]
The first estimate follows by bounding the denominator in
the explicit disk factor (MA-MO.906) below by a fixed positive
constant.  Its constant is independent of the zero and height.

For a cluster in (MA-MO.1203), evaluate at
$z_n=a+d_n+iT_n$, or $w_n=1/2+z_n$ in the uncentered
finite-$L$ disk.  Its distance to every counted centered zero is
at most $5d_n/2$.  Use the first bound for those factors and
the second for the rest.  The deleted finite Euler factor
is uniformly bounded on this positive central strip.  Hence
\[
 |A_\eta(z_n)|\le C\nu(T_n)^D(C_1d_n)^{k_n}.
 \tag{MA-MO.1220}
\]
The usual fixed-window zero bound is $k_n=O(\log\nu(T_n))$,
as in MT-M18 and MA-MO-030.  Thus $C_1^{k_n}$ can be absorbed
in a fixed power of $\nu(T_n)$, and
\[
 d_n^p|A_\eta(a+d_n+iT_n)|^{-1}
           \ge C^{-1}\nu(T_n)^{-D'}d_n^{\,p-k_n}.
 \tag{MA-MO.1221}
\]
Condition (MA-MO.1203) makes the right side exceed every
polynomial in $\nu(T_n)$.  It therefore contradicts
(MA-MO.1201), and the already constructed residual obstruction
applies.  No existence statement about such zeros was used.

Finally the exact identity
$R D_{\eta,g}^0R=D_{\bar\eta,g^\vee}^0$ and the isometry of
$R$ in the sum norm give the reflected version.  The original
finite Fourier phases and primitive Euler conventions remain
those in (MA-MO.550)--(MA-MO.576).
\end{proof}

\paragraph{What this adds, and what it does not decide.}
The new necessity records quantitative information beyond the
individual local lengths in MA-MO.618.  It is a uniform estimate forced by a
whole fixed-support test family on the actual full source.
The cluster implication explains a possible additional analytic
failure mechanism even with only strictly interior, hence fully
retained, local factors.  It does not refute full-retention
sufficiency for actual Hecke arithmetic without an arithmetic
configuration satisfying (MA-MO.1203).  Nor is the reciprocal
estimate asserted sufficient for a trace-class matrix; an
appropriate boundary-distribution estimate is still needed.
MA-MO-030 provides one sufficient class and remains unchanged.

The stipulated input A gives the positive fixed-gap situation for
the Dirichlet family, and M gives it for the stated Hecke family.
Then MA-MO-022 already supplies the full source operators and
(MA-MO.1201) follows.  Invoking that premise does not provide an
independent estimate without zero-location input.  The completed
ideal result MA-MO-028 and the whole-quotient return criteria
MA-PR-009--011 have different domains and are not contradicted.

\paragraph{Source and finite-length placement.}
The original inverse, source membership, dual-term support,
bilinear transpose and compact remainder are those completely
proved in MA-MO-024/025.  Their primary interfaces are Meyer
(2005), supplied \texttt{main.tex} lines 2630--2675,
3592--3707 and 3808--4043, as recorded there.
The fixed-disk analytic ingredients are the finite-$L$ polynomial
strip bound and local zero count used in MT-M18.7--9 and
MA-MO-030; the upper factor estimate (MA-MO.1219)--(MA-MO.1221)
has been derived above.  This module does not claim a new full
reading of those primary proofs.

In the inherited cote~1 trace framework, BR-TR-01/02,
scans 079--080 and 111--112, equality of retained local lengths
is a statement in the category of finite-length
$\mathbb C[z]$-modules.  The height-uniform original-source
bound proved necessary here is additional analytic structure.
Neither that $K_0$ equality, the exterior-power determinants
of the finite underlying vector spaces, nor a motivic weight
supplies it.  No arithmetic progress of FT002 is claimed.
